\subsection{Quartic initialization on the actual cubic Belyi source}

\begin{statement}{A.120}{Proved}
{All quartic common characters vanish on the pointed cubic source}
Put
\[
       \mathcal C=\{0,1,2/3,-1/3\},\qquad
       V=\mathbf P^1\setminus(\mathcal C\cup\{\infty\}).
\tag{A.120.1}
\]
This is the actual source of the degree-three map
$T\mapsto27T^2(1-T)/4$ in A.114.
Let $R,T,c\in\Q_2\setminus\mathcal C$, with $c\ne R,T$,
and suppose that every $(T-a)/(R-a)$, $a\in\mathcal C$,
is a fourth power in $\Q_2^\times$.
For every nonempty subset $J\subset\mathcal C$ and every
$r\in\Q_2^\times$ satisfying
\[
                 \prod_{a\in J}\frac{c-a}{R-a}=-r^2,
\tag{A.120.2}
\]
one has the exact evaluation
\[
                \left(\frac{T-c}{R-c},2r\right)=0
                                \quad\hbox{in }\operatorname{Br}(\Q_2).
\tag{A.120.3}
\]

Consequently, for any rational $R\notin\mathcal C$, any
rational new puncture $c\notin\mathcal C$, any finite set
$S$ of places, and one prescribed ordinary point of
\[
             C_R:\quad T-a=(R-a)X_a^4
                               \quad(a\in\mathcal C)
\tag{A.120.4}
\]
at each $v\in S$, with $T\ne c$, there is
$a_c\in\Q^\times$ such that
\[
 \left\{
 \begin{array}{ll}
 T-a=(R-a)X_a^4 &(a\in\mathcal C),\\
 T-c=a_c Z^4&
 \end{array}\right.
       \longrightarrow C_R,\qquad
       (T,(X_a)_a,Z)\longmapsto(T,(X_a)_a)
\tag{A.120.5}
\]
has an ordinary point over every completion and lifts
every prescribed point.
This is a finite exponent-four initialization statement
on the same five-punctured source as A.114.
\end{statement}
\begin{proof}
Write $v=v_2$ and $f=(T-c)/(R-c)$.
The case $T=R$ is immediate, so assume $T\ne R$.
We first give two local calculations for arbitrary finite
sets of old punctures.

\emph{Equal old valuation coordinates.}
Suppose $v(T-a)=v(R-a)$ at every old puncture $a$.
A unit of $\Q_2$ is a fourth power exactly when it lies
in $1+16\Z_2$; this follows by twice applying the square
criterion, or from the isomorphism given by the two-adic
logarithm on $1+4\Z_2$.
Thus, with $M=\max_a v(R-a)$ and $s=v(T-R)$,
\[
       \frac{T-a}{R-a}=1+\frac{T-R}{R-a}\in1+16\Z_2,
                           \qquad s\ge M+4.
\tag{A.120.6}
\]
If $v(c-R)\ge M+2$, every $(c-a)/(R-a)$ is in
$1+4\Z_2$.  No product of these units is a negative
square, since a negative unit square is $3$ modulo four.
Hence (A.120.2) forces $v(c-R)\le M+1$.
Equation (A.120.6) then gives $f\in1+8\Z_2$, so $f$ is
a square and (A.120.3) follows.

\emph{Two points deep in one boundary disc.}
Consider any old set consisting of $0$ and nonzero numbers
$b$ with $v(b)\le0$.
Assume
\[
             e=v(R)\ge4,\quad e_T=v(T)\ge4,
             \quad T/R=X^4,\quad e_T\equiv e\pmod4.
\tag{A.120.7}
\]
We prove (A.120.3) for every negative subset product in
this set; the number of other punctures is arbitrary.
For $k=v(c)\le0$, $f\in1+16\Z_2$, so it suffices to
consider $k>0$.

If the subset does not contain zero, its product is
$G(c)/G(R)$, where
$G(z)=\prod_b(1-z/b)$ over that subset.
For $k\ge2$ this product is in $1+4\Z_2$ and cannot be
a negative square.  For $k=1$, one has $f\in1+8\Z_2$.
This handles every such subset.

For a subset containing zero, use the same $G$ for its
remaining punctures; the empty product is one.
The exact relation is
\[
                    r^2=-\frac cR\frac{G(c)}{G(R)}.
\tag{A.120.8}
\]
When $k\ge4$, the quotient $G(c)/G(R)$ lies in
$1+16\Z_2$.  Choose its square root $l\in1+8\Z_2$;
the number $l$ is itself a square.
Put $a=r/l$, so $a^2=-c/R$ and $2a$ has the same
quadratic character as $2r$.
The literal norm identity is
\[
 \begin{split}
 f&=\frac{X^4+a^2}{1+a^2},\\
 X^4+a^2&=(X^2+a)^2-2aX^2,
       \qquad 1+a^2=(1+a)^2-2a.
 \end{split}
\tag{A.120.9}
\]
It proves $(f,2r)=(f,2a)=0$.

It remains to treat $k=1,2,3$.
The factors of $G(c)$ and $G(R)$ are units, so
(A.120.8) gives $k\equiv e\pmod2$ and
$v(2r)=1+(k-e)/2$.
Put $m=\min(e,e_T)$.
Both $e$ and $e_T$ are congruent modulo four; hence
$m-k$ is a positive even integer.
If $v(T-R)-k\ge3$, then $f$ is a square.
The remaining possibility is
\[
 m-k=2,\qquad v(T-R)=m,\qquad
 e\ne e_T,\qquad e\equiv m\pmod4,
                       \qquad v(2r)\equiv0\pmod2.
\tag{A.120.10}
\]
Here the inequality $e\ne e_T$ is forced: for equal
depths, $T/R=X^4$ would give $v(T-R)\ge e+4$.
In (A.120.10), $f\equiv5\pmod8$.
The extension $\Q_2(\sqrt f)=\Q_2(\sqrt5)$ is unramified
quadratic, with norm group the elements of even valuation.
Thus $2r$ is a norm, proving the deep-disc calculation.

\emph{The exact valuation alternatives on (A.120.1).}
In the order $0,1,2/3,-1/3$, the valuation vector of an
ordinary $z\in\Q_2$ has exactly one of the following forms:
\[
\begin{array}{c|c|c}
\text{location}&(v(z),v(z-1),v(z-2/3),v(z+1/3))
                                      &\text{range}\\ \hline
v(z)<0&(h,h,h,h)&h<0\\
\text{near }0&(e,0,1,0)&e\ge2\\
\text{near }2/3&(1,0,e,0)&e\ge2\\
\text{odd middle disc}&(0,1,0,1)&\\
\text{near }1&(0,e,0,2)&e\ge3\\
\text{near }-1/3&(0,2,0,e)&e\ge3.
\end{array}
\tag{A.120.11}
\]
Indeed $v(2/3)=1$ and $v(1+1/3)=2$.
For even $z$, exactly one of $z,z-2/3$ has valuation one
and the other has valuation at least two.
For odd $z$, the last two relevant distances are either
both one, or one is two and the other is at least three.
These observations prove every row and exhaust $\Q_2$.

The fourth-power ratios make the two valuation vectors
componentwise congruent modulo four.
The negative row cannot be congruent to an integral row,
which always contains a zero and either a one or a two.
The two even rows cannot be congruent to an odd row:
one of their first and third entries is one, whereas
both corresponding entries of every odd row are zero.
The odd middle row cannot be congruent to either deeper
odd row, because one of the latter's two nonzero entries
is two.

A change from the $0$ row to the $2/3$ row would require
both depths to be $1$ modulo four, hence at least five.
After interchanging the names of the two points in this
contradiction argument, $R\equiv0$ and $T\equiv2/3$
modulo $32$.
Consequently
\[
       \frac{T-1}{R-1}\equiv\frac13\equiv11\pmod{16},
\tag{A.120.12}
\]
which is not a fourth-power unit.
A change between the $1$ and $-1/3$ rows would require
both depths to be $2$ modulo four, hence at least six.
The ratio at zero would then satisfy
\[
                         T/R\equiv-1/3\equiv5\pmod{16},
\tag{A.120.13}
\]
again impossible.
Thus either all four valuations agree, or the two points
have different depths in the same boundary row of
(A.120.11).

\emph{Finite boundary rows and their actual affine maps.}
Suppose the common row is near $a\in\mathcal C$, with
distinct depths $e,e_T$.
Let
$\delta=\max_{b\in\mathcal C\setminus\{a\}}v(a-b)$.
It equals one for $a=0,2/3$, and two for $a=1,-1/3$.
For a neighbour attaining this maximum, both distances
$R-b,T-b$ have valuation $\delta$.
The fourth-power unit ratio at this neighbour gives
\[
                \min(e,e_T)=v(T-R)\ge\delta+4.
\tag{A.120.14}
\]
Use the affine map $z\mapsto(z-a)/2^\delta$ on
$R,T,c$ and all the punctures.
Every difference ratio, the normalized entry $f$, and
the subset product in (A.120.2) is literally unchanged.
On (A.120.4) this map divides each coefficient by
$2^\delta$ and keeps the root coordinates $X_b$.
The distinguished old puncture becomes zero, every
other old puncture has valuation at most zero, and
(A.120.14) gives both new depths at least four.
The deep-disc calculation therefore applies.

\emph{The infinite row and its exact root-coordinate map.}
Finally suppose $v(R),v(T)<0$ are distinct, and put
$e=-v(R)$, $e_T=-v(T)$.
Dividing the fourth-power ratio at one by the one at
zero gives
\[
       \frac{1-1/T}{1-1/R}\in\Q_2^{\times4},
       \qquad
       v\left(\frac{1-1/T}{1-1/R}-1\right)
                          =\min(e,e_T).
\tag{A.120.15}
\]
This is a unit, so both depths are at least four.
Apply $z\mapsto1/z$ to the full boundary set.
Its new finite punctures are
$0,1,3/2,-3$, whose nonzero members have valuations
$0,-1,0$.
The explicit map of pointed quartic covers is
\[
 U=1/T,\qquad Y_0=X_0^{-1},\qquad
 Y_{1/b}=X_b/X_0\quad(b=1,2/3,-1/3).
\tag{A.120.16}
\]
Its inverse is $T=1/U$, $X_0=Y_0^{-1}$ and
$X_b=Y_{1/b}/Y_0$.
It identifies the old point and boundary-fibre algebras;
the new difference ratios are $R/T$ and
$(R/T)(T-b)/(R-b)$, respectively.

For completeness the negative-subset characters can
also be transported without any common-field assertion
over a local field.
Write $D_0=c/R$, $D_b=(c-b)/(R-b)$, and let primes
denote the new data.
The actual equalities are
\[
 D'_0=D_0^{-1},\qquad D'_{1/b}=D_b/D_0,
            \qquad f'=fR/T.
\tag{A.120.17}
\]
For a subset $J$ of size $j$, form $J'$ from its nonzero
members under $b\mapsto1/b$, adjoining zero exactly
when $j$ is odd.
Then
\[
 \prod_{b'\in J'}D'_{b'}
       =-\bigl(rD_0^{-\lceil j/2\rceil}\bigr)^2,
       \qquad r'=rD_0^{-\lceil j/2\rceil}.
\tag{A.120.18}
\]
The individual norm identity (A.118.13), for the old
puncture zero, proves $(f,D_0)=0$.
Since $R/T$ is a fourth power, (A.120.17)--(A.120.18)
give
\[
          (f',2r')=(f,2r)
                      \quad\hbox{in }\operatorname{Br}(\Q_2).
\tag{A.120.19}
\]
Both new points have depth at least four at zero.
The deep-disc calculation proves that the left side is
zero.  This finishes every case of (A.120.3).

\emph{The full finite-character initialization on this source.}
For rational $R,c$ with $c\ne R$, A.96 describes every
common character of the boundary fibre of (A.120.4):
the individual quadratic characters and all characters
of $2r_J$ from rational negative-square subset products.
The individual normalized classes vanish by their exact
quadratic norms (A.118.13).
All the negative-subset classes vanish at two by the
proof above, and at every odd and real place by A.118.
Thus every character has zero normalized local
evaluation at every old local point avoiding $c$.
Subtracting the global constant $(R-c,\chi)$ makes the
sum in the exact criterion A.85 zero by reciprocity.
That criterion supplies the rational coefficient $a_c$
and the actual forgetful map (A.120.5), retaining the
finitely prescribed points.
For $c=R$, the boundary fibre has the rational root
tuple $(1,1,1,1)$, so its common field is $\Q$ and the
same criterion has no nonzero characters to test.
This completes the finite initialization assertion.
No compatibility between coefficients chosen at
successive exponents is used or asserted.
\end{proof}

The exact finite checks are retained in
\path{code/cubic_source_quartic_certificate.py} and its
output.  They audit the displayed cover and subset maps,
the valuation alternatives and the local Hilbert
evaluations; the all-point assertion is proved above.
